\documentclass[aspectratio=169]{beamer}
\usepackage{etex} % fixes new-dimension error
\usepackage{lmodern}
\usepackage[T1]{fontenc}

\usepackage{graphicx,amsmath}
\usepackage{stmaryrd} % cf. interleave
\usepackage{booktabs}
\usepackage{amscd}
\usepackage{multicol}
\usepackage[absolute,overlay]{textpos}
\usepackage{alltt}
\usepackage{proof}
\usepackage{subcaption}
\usepackage{tikz}
\usepackage{tikz-cd}
\usepackage[new]{old-arrows}
\usepackage[all]{xy}
\usepackage{pgfplots}
\usepackage{textcomp}

\usepackage{transparent}
\usepackage{xspace}
\usepackage{listings}
\usepackage{pdfpages}
\usepackage{relsize}

%%%%%%%%%%%%% Macros
\newcommand{\Ban}{\catfont{Ban}}
\newcommand{\Met}{\catfont{Met}}
\newcommand{\Shuff}{\mathrm{Sf}}
\newcommand{\Cats}{\catfont{Cat}}
\newcommand{\VCat}{\mathcal{V}\text{-}\Cats}
\newcommand{\VCatSy}{\mathcal{V}\text{-}\Cats_{\mathsf{sym}}}
\newcommand{\VCatSe}{\mathcal{V}\text{-}\Cats_{\mathsf{sep}}}
\newcommand{\VCatSS}{\mathcal{V}\text{-}\Cats_{\mathsf{sym,sep}}}
%%%% Categories
\newcommand{\catfont}[1]{\mathsf{#1}}
\newcommand{\cop}{\catfont{op}}
\newcommand{\Law}{\catfont{Law}}
\newcommand{\catV}{\catfont{V}}
\newcommand{\catX}{\catfont{X}}
\newcommand{\catC}{\catfont{C}}
\newcommand{\catD}{\catfont{D}}
\newcommand{\catA}{\catfont{A}}
\newcommand{\catB}{\catfont{B}}
\newcommand{\catI}{\catfont{I}}
\newcommand{\Set}{\catfont{Set}}
\newcommand{\Top}{\catfont{Top}}
\newcommand{\Pos}{\catfont{Pos}}
\newcommand{\Inj}{\catfont{Inj}}
\newcommand{\Det}{\catfont{RMhat}}
\newcommand{\CoAlg}[1]{\catfont{CoAlg}\left (#1 \right )}
\newcommand{\Mon}{\catfont{Mon}}
\newcommand{\Mnd}{\catfont{Mnd}(\catC)}
\newcommand{\SMnd}{\catfont{Mnd}(\Set)}
\newcommand{\CLat}{\catfont{CLat}}
\newcommand{\Stone}{\catfont{Stone}}
\newcommand{\Spectral}{\catfont{Spectral}}
\newcommand{\CompHaus}{\catfont{CompHaus}}
\newcommand{\Subs}[2]{\catfont{Sub}_{}}
\newcommand{\Cone}{\catfont{Cone}}
\newcommand{\StComp}{\catfont{StablyComp}}
\newcommand{\PosC}{\catfont{PosComp}}
\newcommand{\Haus}{\catfont{Haus}}
\newcommand{\Meas}{\catfont{Meas}}
\newcommand{\Ord}{\catfont{Ord}}
\newcommand{\EndoC}{[\catC,\catC]}
%% General functors
\newcommand{\funfont}[1]{#1}
\newcommand{\funF}{\funfont{F}}
\newcommand{\funU}{\funfont{U}}
\newcommand{\funG}{\funfont{G}}
\newcommand{\funT}{\funfont{T}}
\newcommand{\funI}{\funfont{I}}
%% Particular kinds of functors
\newcommand{\sfunfont}[1]{\mathrm{#1}}
\newcommand{\Pow}{\sfunfont{P}}
\newcommand{\Dist}{\sfunfont{D}}
\newcommand{\Maybe}{\sfunfont{M}}
\newcommand{\List}{\sfunfont{L}}
\newcommand{\UForg}{\sfunfont{U}}
\newcommand{\Forg}[1]{\sfunfont{U}_{#1}}
\newcommand{\Id}{\sfunfont{Id}}
\newcommand{\Vie}{\sfunfont{V}}
\newcommand{\Disc}{\funfont{D}}
\newcommand{\Weight}{\sfunfont{W}}
\newcommand{\homf}{\sfunfont{hom}}
\newcommand{\Yoneda}{\sfunfont{Y}}
%% Diagram functors
\newcommand{\Diag}{\mathscr{D}}
\newcommand{\KDiag}{\mathscr{K}}
\newcommand{\LDiag}{\mathscr{L}}
%% Monads
\newcommand{\monadfont}[1]{\mathbb{#1}}
\newcommand{\monadT}{\monadfont{T}}
\newcommand{\monadS}{\monadfont{S}}
\newcommand{\monadU}{\monadfont{U}}
\newcommand{\monadH}{\monadfont{H}}
\newcommand{\str}{\mathrm{str}}
%% Adjunctions
\newcommand\adjunct[2]{\xymatrix@=8ex{\ar@{}[r]|{\top}\ar@<1mm>@/^2mm/[r]^{{#2}}
& \ar@<1mm>@/^2mm/[l]^{{#1}}}}
\newcommand\adjunctop[2]{\xymatrix@=8ex{\ar@{}[r]|{\bot}\ar@<1mm>@/^2mm/[r]^{{#2}}
& \ar@<1mm>@/^2mm/[l]^{{#1}}}}
%% Retractions
\newcommand\retract[2]{\xymatrix@=8ex{\ar@{}[r]|{}\ar@<1mm>@/^2mm/@{^{(}->}[r]^{{#2}}
& \ar@<1mm>@/^2mm/@{->>}[l]^{{#1}}}}
%% Limits
\newcommand{\pv}[2]{\langle #1, #2 \rangle}
\newcommand{\limt}{\mathrm{lim}}
\newcommand{\pullbackcorner}[1][dr]{\save*!/#1+1.2pc/#1:(1,-1)@^{|-}\restore}
\newcommand{\pushoutcorner}[1][dr]{\save*!/#1-1.2pc/#1:(-1,1)@^{|-}\restore}
%% Colimits
\newcommand{\colim}{\mathrm{colim}}
\newcommand{\inl}{\mathrm{inl}}
\newcommand{\inr}{\mathrm{inr}}
%% Distributive categories
\newcommand{\distr}{\mathrm{dist}}
\newcommand{\undistr}{\mathrm{undist}}
%% Closedness
\newcommand{\curry}[1]{\mathrm{curry}{#1}}
\newcommand{\app}{\mathrm{app}}
%% Misc. operations
\newcommand{\const}[1]{\underline{#1}}
\newcommand{\comp}{\cdot}
\newcommand{\id}{\mathrm{id}}
%% Factorisations
\newcommand{\EClass}{E}
\newcommand{\MClass}{M}
\newcommand{\MConeClass}{\mathcal{M}}
%%%%%%%%%%%%%%%% End of Categorical Stuff

%%%% Misc
%% Operations
\newcommand{\blank}{\, - \,}
\newcommand{\sem}[1]{\llbracket #1 \rrbracket}
\newcommand{\closure}[1]{\overline{#1}}
\DeclareMathOperator{\img}{\mathrm{im}}
\DeclareMathOperator{\dom}{\mathrm{dom}}
\DeclareMathOperator{\codom}{\mathrm{codom}}
%% Sets of numbers
\newcommand{\N}{\mathbb{N}}
\newcommand{\Z}{\mathbb{Z}}
\newcommand{\Nats}{\mathbb{N}}
\newcommand{\Reals}{\mathbb{R}}
\newcommand{\Rz}{\Reals_{\geq 0}}
\newcommand{\Complex}{\mathbb{C}}
%% Writing
\newcommand{\cf}{\emph{cf.}}
\newcommand{\ie}{\emph{i.e.}}
\newcommand{\eg}{\emph{e.g.}}
\newcommand{\df}[1]{\emph{\textbf{#1}}}
%%%%%%%%%%%%%%%% End of Misc

%%%% Programming Stuff
%% Types
\newcommand{\typefont}[1]{\mathbb{#1}}
\newcommand{\typeOne}{1}
\newcommand{\typeTwo}{2}
\newcommand{\typeA}{\typefont{A}}
\newcommand{\typeX}{\typefont{X}}
\newcommand{\typeB}{\typefont{B}}
\newcommand{\typeC}{\typefont{C}}
\newcommand{\typeV}{\typefont{V}}
\newcommand{\typeD}{\typefont{D}}
\newcommand{\typeI}{\typefont{I}}
%% RuleName
\newcommand{\rulename}[1]{(\mathrm{#1})}
%% Sequents
\newcommand{\jud}{\vdash}
\newcommand{\vljud}{\rhd}
\newcommand{\cojud}{\vdash_{\co}}
\newcommand{\vl}{\mathtt{v}}
\newcommand{\co}{\mathtt{c}}
% Program font
\newcommand{\prog}[1]{\mathtt{#1}}
\newcommand{\pseq}[3]{#1 \leftarrow #2; #3}
\newcommand{\ppm}[4]{(#1,#2) \leftarrow #3; #4}
\newcommand{\pinl}[1]{\prog{inl}(#1)}
\newcommand{\pinr}[1]{\prog{inr}(#1)}
\newcommand{\pcase}[4]{\prog{ case } #1 \prog{ of } \pinl{#2} \Rightarrow #3 ; \pinr{#2} \Rightarrow #4}
%% Sets of terms
\newcommand{\ValuesBP}[2]{\mathsf{Values}(#1, #2)}
\newcommand{\TermsBP}[2]{\mathsf{Terms}(#1, #2)}
\newcommand{\closValP}[1]{\ValuesBP{\emptyset}{#1}}
\newcommand{\closTermP}[1]{\TermsBP{\emptyset}{#1}}
\newcommand{\closVal}{\closValP{\typeA}}
\newcommand{\closTerm}{\closTermP{\typeA}}
%% Contextual equivalence
\newcommand{\ctxeq}{\equiv_{\prog{ctx}}}
%%%% End of Programming Stuff

%------ Setting lecture info ----------------------------------------------
\newcounter{lectureID}
\stepcounter{lectureID}
\newcommand{\getLecture}{\arabic{lectureID}\xspace}
\newcommand{\setLectureBasic}[1]{
  \title{
    #1
    }
  \author{Jos\'{e} Proen\c{c}a}
  \institute{CISTER -- U.Porto, Porto, Portugal
            \hfill 
            \begin{tabular}{r@{}}
            \url{https://fm-dcc.github.io/alg2425}
            \end{tabular}
            }
  \date{Algorithms (CC4010) 2024/2025}
  % logos of institutions
  \titlegraphic{
    \begin{textblock*}{5cm}(4.0cm,6.80cm)
       \includegraphics[scale=0.18]{images/fcup}\hspace*{.85cm}~%
    \end{textblock*}
    \begin{textblock*}{5cm}(8.4cm,7.25cm)
      % \includegraphics[scale=0.50]{images/dcc}
      \includegraphics[scale=0.20]{images/cister}
    \end{textblock*}
  }  
}
\newcommand{\setLecture}[2]{\setcounter{lectureID}{#1}\setLectureBasic{#1. #2}}

%------ Counters for exercises ----------------------------------------------
\newcounter{cExercise}
\newcommand{\exercise}{\stepcounter{cExercise}Ex.\,\arabic{lectureID}.\arabic{cExercise}:\xspace}
\newcommand{\exerciseBack}{\addtocounter{cExercise}{-1}}
\newcommand{\exerciseAdd}{\stepcounter{cExercise}}
\newcommand{\doExercise}[3][0mm]{\begin{exampleblock}{\exercise #2}\wrap{\rule{0pt}{#1}}#3\end{exampleblock}}
\newcommand{\doSimpleExercise}[2][0mm]{\begin{exampleblock}{}\wrap{\rule{0pt}{#1}}\structure{\textbf{\exercise} #2}\end{exampleblock}}

% Slide
\newenvironment{slide}[1]{\begin{frame}\frametitle{#1}}{\end{frame}}

% Misc by José
\newcommand{\wrap}[2][]{\begin{tabular}[#1]{@{}c@{}}#2\end{tabular}}
\newcommand{\mwrap}[1]{\ensuremath{\begin{array}{@{}c@{}}#1\end{array}}}
\def\trans#1{\xrightarrow{#1}}  % - a - > 
\def\Trans#1{\stackrel{#1}{\Longrightarrow}} % =a=> 
\newcommand{\transp}[2][35]{\color{fg!#1}#2}
\newcommand{\transpt}[2][.35]{\tikz{\node[inner sep=1pt,fill opacity=0.5]{#2}}}
\newcommand{\faded}[2][0.4]{{\transparent{#1}#2}} % alternative to "transp" using transparent package
\newcommand{\set}[1]{\left\{ #1 \right\}} % {a,b,...z}
\newcommand{\mi}[1]{\ensuremath{\mathit{#1}}\xspace}
\newcommand{\mf}[1]{\ensuremath{\mathsf{#1}}\xspace}
% \newcommand{\gold}[1]{\textcolor{darkgoldenrod}{#1}\xspace}


%------ using color ---------------------------------------------------------
\definecolor{goldenrod}{rgb}{.80392 .60784 .11373}
\definecolor{darkgoldenrod}{rgb}{.5451 .39608 .03137}
\definecolor{brown}{rgb}{.15 .15 .15}
\definecolor{darkolivegreen}{rgb}{.33333 .41961 .18431}
\definecolor{myGray}{gray}{0.85}
%
%
\newcommand{\red}[1]{\textcolor{red!80!black}{#1}\xspace}
\newcommand{\blue}[1]{\textcolor{blue}{#1}\xspace}
\newcommand{\gold}[1]{\textcolor{darkgoldenrod}{#1}\xspace}
\newcommand{\gray}[1]{\textcolor{myGray}{#1}\xspace}
% \def\alert#1{{\darkgoldenrod #1}}
% \def\alert#1{{\alert{#1}}}
%\def\brw#1{{\brown #1}}
% \def\structure#1{{\blue #1}}
% \def\tstructure#1{\textbf{\darkblue #1}}
%%\def\gre#1{{\green #1}}
\def\gre#1{{\darkolivegreen #1}}
\def\gry#1{{\textcolor{gray}{#1}}}
\def\rdb#1{{\red #1}}
\def\st{\mathbf{.}\,}
\def\laplace#1#2{*\txt{\mbox{ \fcolorbox{black}{myGray}{$\begin{array}{c}\mbox{#1}\\\\#2\\\\\end{array}$} }}}
%\newcommand{\galois}[2]{#1\; \dashv\; #2}



% ----- from LSB
\def\Act{N}
\def\cnil{\mathbf{0}}
\def\cpf#1#2{#1 . #2}                           % a.P
\def\cou#1#2{#1 \mathbin{+} #2}                 % P + Q
%\def\crt#1#2{\mathbin{#1 \setminus_{#2}}}       % P \ A
%\def\crtt#1#2{\mathbin{#1 \setminus\!\setminus_{#2}}}       % P \ A
%\def\crt#1#2{\mathsf{new}\, #2\;  #1}       % P \ A
\def\crt#1#2{#1 \backslash #2}       % P \ A
%\def\crn#1#2{\{#2\}\, #1}                  % P[f]
\def\ainv#1{\overline{#1}}
\def\rtran#1{\stackrel{#1}{\longrightarrow}}
% \def\pair#1{\const{#1}}
\def\pair#1{\langle #1 \rangle}
\def\asor{\mathbin{|}}                    % A | B
\def\setdef#1#2{\mathopen{\{} #1 \asor #2 \mathclose{\}}}
\def\imp{\mathbin{\Rightarrow}}
\def\dimp{\mathbin{\Leftrightarrow}}
\def\rimp{\mathbin{\Leftarrow}}
\def\rra{\longrightarrow}
\def\rcb#1#2#3#4{\def\nothing{}\def\range{#3}\mathopen{\langle}#1 \ #2 \ \ifx\range\nothing::\else: \ #3 :\fi \ #4\mathclose{\rangle}}
\def\aconv#1{#1^{\circ}} 
\def\abv{\stackrel{\rm abv}{=}}


\def\crn#1#2{\mathbin{#1[#2]}}                  % P[f]
\def\couit#1#2{\Sigma_{#1}#2}                  %  + i=1,n
\def\cpar#1#2{#1 \mid #2}                       %  |
\def\ctpar#1#2{#1 \parallel #2}                       %  |
\def\cpars#1#2#3{#1 \mid_{#3} #2}               %  |S

% Spliting frames in 2 columns
\newcommand{\frsplit}[3][.48]{
  \begin{columns}%[T] % align columns
  \begin{column}{#1\textwidth} #2 \end{column} ~~~
  \begin{column}{#1\textwidth} #3 \end{column} \end{columns}
}
\newcommand{\frsplitdiff}[5][]{
  \begin{columns}[#1]%[T] % align columns
  \begin{column}{#2\textwidth} #4 \end{column} ~~~
  \begin{column}{#3\textwidth} #5 \end{column} \end{columns}
}
\newcommand{\frsplitt}[3][.48]{
  \begin{columns}[T] % align columns
  \begin{column}{#1\textwidth} #2 \end{column} ~~~
  \begin{column}{#1\textwidth} #3 \end{column} \end{columns}
}
\newcommand{\col}[2][.48]{\begin{column}{#1\textwidth} #2 \end{column}}
\newcommand{\colb}[3][.48]{\begin{column}{#1\textwidth} \begin{block}{#2} #3 \end{block} \end{column}}

% Spliting frames in 3 columns
\newcommand{\frsplitthree}[4][.31]{
  \begin{columns}%[T] % align columns
  \begin{column}{#1\textwidth} #2 \end{column} ~~~
  \begin{column}{#1\textwidth} #3 \end{column} ~~~
  \begin{column}{#1\textwidth} #4 \end{column} \end{columns}
}
\newcommand{\frsplitdiffthree}[5]{
  \begin{columns}%[T] % align columns
  \begin{column}{#1\textwidth} #3 \end{column} ~~~
  \begin{column}{#1\textwidth} #4 \end{column} ~~~
  \begin{column}{#2\textwidth} #5 \end{column} \end{columns}
}
\newcommand{\frsplittthree}[4][.32]{
  \begin{columns}[T] % align columns
  \begin{column}{#1\textwidth} #2 \end{column} ~
  \begin{column}{#1\textwidth} #3 \end{column} ~
  \begin{column}{#1\textwidth} #4 \end{column} \end{columns}
}


\newcommand{\typerule}[4][]{\ensuremath{\begin{array}[#1]{c}\textsf{\scriptsize ({#2})} \\#3 \\\hline\raisebox{-3pt}{\ensuremath{#4}}\end{array}}}
\newcommand{\styperule}[3][]{\ensuremath{\begin{array}[#1]{c} #2 \\[0.5mm]\hline\raisebox{-4pt}{\ensuremath{#3}}\end{array}}}
\newcommand{\shrk}{\vspace{-3mm}}

\def\caixa#1{\medskip
  \begin{center}
  \fbox{\begin{minipage}{0.9\textwidth}\protect{#1}\end{minipage}}
  \end{center}}

\newcommand{\mybox}[2][0.9]{
  \begin{minipage}{#1\textwidth}\begin{block}{}\centering #2\end{block}\end{minipage}}
\newcommand{\mycbox}[2][0.9]{
  {\\[-5mm]\centering\mybox[#1]{#2}\\[-5mm]}}

%%%%% Tikz
% \usetikzlibrary{arrows.meta, calc, fit, tikzmark}
\usetikzlibrary{%
  positioning
 ,patterns
 ,arrows
 ,arrows.meta
 ,automata
 ,calc
 ,shapes
 ,fit
 ,tikzmark
 ,fadings
 ,decorations.pathreplacing
 ,plotmarks
% ,pgfplots.groupplots
 ,decorations.markings
}
% \tikzset{shorten >=1pt,node distance=2cm,on grid,auto,initial text={},inner sep=2pt}
\tikzstyle{aut}=[shorten >=1pt,node distance=2cm,on grid,auto,initial text={},inner sep=2pt]
\tikzstyle{st}=[circle,draw=black,fill=black!10,inner sep=3pt]
\tikzstyle{sst}=[rectangle,draw=none,fill=none,inner sep=3pt]
\tikzstyle{final}=[accepting]

%%% Uppaal-like diagrams
\newcommand{\uppbox}[3][20mm]{\tikz{
  \node[black!15,fill=black!15,minimum width=#1,align=left](title){\textbf{{\footnotesize #2}}};
  \node[black!15,fill=black!15,left,xshift=4mm]at(title.east){\textbf{{\footnotesize #2}}};
  \node[blue!60!cyan,right] at(title.west){\textbf{{\footnotesize #2}}};
  \node[below,inner sep=2mm,fill=white,xshift=2mm](box)at(title.south){\includegraphics[width=#1]{#3}};
  \node[fit=(title)(box),draw=black,inner sep=0pt]{};
}}

\newcommand{\uppboxv}[3][20mm]{\tikz{
  \node[below,inner sep=2mm,fill=white](box){\includegraphics[height=#1]{#3}};
  \coordinate[yshift=5mm](top)at(box.north);
  \node[fit=(top)(box.north west)(box.north east),inner sep=0pt,fill=black!15](title){};
  \node[blue!60!cyan,right] at(title.west){\textbf{{\footnotesize #2}}};
  \node[fit=(title)(box),draw=black,inner sep=0pt]{};
}}


%% COnfiguring Listings
\lstset{ % basic style
  % language=scala,
  basicstyle=\ttfamily\scriptsize,
  breakatwhitespace=true,
  breaklines=true,
  mathescape,
  % morecomment=[l]{//},
  % morecomment=[n]{/*}{*/},
  % frame=single,                    % adds a frame around the code
  rulecolor=\color{black!40},         % if not set, the frame-color may be changed on line-breaks within not-black text (e.g. comments (green here))
  xleftmargin=1.5mm,
  xrightmargin=1.5mm,
  backgroundcolor=\color{black!5},
  % line numbers
%  numbers=left,  % where to put the line-numbers; possible values are (none, left, right)
 numbersep=5pt, % how far the line-numbers are from the code
 numberstyle=\tiny\color{gray},   
 stepnumber=1,  % the step between two line-numbers. If it is 1 each line will be numbered      
%  xleftmargin=3mm,
%  xrightmargin=1.5mm,
%%%%%
  captionpos=b, % t or b (top or bottom)
  belowcaptionskip=5mm,
%%%%%
  % alsoletter={-},
  % emphstyle=\ttfamily\color{blue}, %\underbar,
  % emphstyle={[2]\ttfamily\color{green!50!black}},
  emphstyle=\bfseries\itshape\color{blue!80!black},       % moreemph={...} - layer keywords
  emphstyle={[2]\itshape\color{red!70!black}},%\underbar} % moreemph={[2]...} - inner keywords
  %
  keywordstyle=\bf\ttfamily\color{red!50!black},
%  commentstyle=\sl\ttfamily\color{gray!70},
  commentstyle=\color{green!60!black},
  stringstyle=\ttfamily\color{purple!60!black},
  morestring=[b]",
  morecomment=[l]{\#},
  frame=single,
  % numberstyle=\tiny, numbers=left, stepnumber=1, firstnumber=1, numberfirstline=true,
  %emph={act,proc,init,sort,map,var,eqn},
  %emph={[2]block,hide,comm,rename,allow,||,<>,sum,&&,=>,true,false},
  % literate=*{->}{{{\color{red!70!black}$\to$}}}{1}
  %            {.}{{{\color{red!70!black}.}}}{1}
  %            {+}{{{\color{red!70!black}\hspace*{1pt}+\hspace*{1pt}}}}{1}
  %            {|}{{{\color{red!70!black}|}}}{1}
  %            {||}{{{\color{red!70!black}|\!\!|}}}{1}
  %            {*}{{{\color{red!70!black}*}}}{1}
  %            {\#}{{{\color{red!70!black}\#}}}{1}
  %            {&}{{{\color{red!70!black}\&}}}{1}
  %            {:=}{{{\color{red!70!black}:=}}}{1}
  %            {=>}{{{\color{red!70!black}=>}}}{1}
  %            % {>}{{{\color{green!65!black}\hspace*{1pt}>\hspace*{1.5pt}}}}{2}
  %            % {<}{{{\color{green!65!black}\hspace*{1.5pt}<\hspace*{1pt}}}}{2}
  %            % {]}{{{\color{green!65!black}\hspace*{1pt}]\hspace*{1.5pt}}}}{1}
  %            % {[}{{{\color{green!65!black}\hspace*{1.5pt}[\hspace*{1pt}}}}{1}
  %            {>}{{{\color{green!65!black}>}}}{1}
  %            {<}{{{\color{green!65!black}<}}}{1}
  %            {]}{{{\color{green!65!black}]}}}{1}
  %            {[}{{{\color{green!65!black}[}}}{1}
  %  morekeywords={Merger1,Fifo2,Lossy3,Init1,Init2}
  % ,emph={act,proc,init,sort,map,var,eqn}
  % ,emph={[2]block,hide,comm,rename,allow,||,<>,sum,&&,=>}
}

\lstdefinestyle{tiny}{basicstyle=\ttfamily\relsize{-7}}

% Include slides from others
\newcommand{\byothers}[3]{{
\begin{frame}{}~\mycbox{\Large slides by #1\\pages #2}\end{frame}{}
\setbeamercolor{background canvas}{bg=}
\includepdf[pages=#2]{../../others/#3}
}}

%-------------- template --------------------------------------------------
\usetheme{metropolis}
\usepackage{appendixnumberbeamer}

% Base colors (from metropolis theme)
\definecolor{metDarkBrown}{HTML}{604c38}
\definecolor{metDarkTeal}{HTML}{23373b}
\definecolor{metLightBrown}{HTML}{EB811B}
\definecolor{metLightGreen}{HTML}{14B03D}

 

\metroset{numbering=fraction,progressbar=frametitle}

% \setbeamercolor*{structure}{fg=blue!80!black}
\setbeamercolor*{structure}{fg=metLightGreen}

% % \definecolor{MainColour}{rgb}{0., 0.25, 0.8}
% \colorlet{MainColour}{blue!50!black}
% \colorlet{BgColour}{blue!10}
% \colorlet{BarColour}{blue!50!black}


% %\usetheme{CambridgeUS}%{Copenhagen}%{Frankfurt}%{Singapore}%{CambridgeUS}
% \usecolortheme[named=MainColour]{structure} 
% \useoutertheme[subsection=false]{miniframes}
% \useinnertheme{circles}
% %\useinnertheme[shadow=false]{rounded}
% \setbeamertemplate{blocks}[rounded][shadow=false]

% \setbeamercovered{transparent} 
% \setbeamertemplate{navigation symbols}{} %Remove navigation bar
% \setbeamertemplate{footline}[frame number] % add page number
% \setbeamercolor{postit}{fg=MainColour,bg=BgColour}
% \setbeamercolor{structure}{bg=black!10}
% %\setbeamercolor{palette primary}{use=structure,fg=red,bg=green}
% %\setbeamercolor{palette secondary}{use=structure,fg=red!75!black,bg=green}
% \setbeamercolor{palette tertiary}{use=structure,bg=BarColour,fg=white}
% %\setbeamercolor{palette quaternary}{fg=black,bg=green}
% %\setbeamercolor{normal text}{fg=black,bg=white}
% %\setbeamercolor{block title alerted}{fg=red,bg=green}
% %\setbeamercolor{block title example}{bg=black!10,fg=green}
\setbeamercolor{block body}{bg=black!5}

% \setbeamercolor{block title alerted}{bg=red!25}
% \setbeamercolor{block body alerted}{bg=red!10}

% \setbeamercolor{block title example}{bg={rgb:green,2;black,1;white,5}}
\setbeamercolor{block body example}{bg={rgb:green,2;black,1;white,20}}
\setbeamercolor{block body alerted}{bg={metLightBrown!25}}

% \setbeamertemplate{itemize item}{\color{black!10}$\blacksquare$}
\setbeamercolor{itemize item}{fg=metDarkTeal}
\setbeamercolor{itemize subitem}{fg=metDarkTeal}

\setbeamercolor{graybc}{fg=black,bg=black!10}
\newcommand{\myblock}[1]{\begin{beamercolorbox}[dp=1ex,center,rounded=true]%
  {graybc} {\large \textbf{#1}} \end{beamercolorbox}}%



\definecolor{barcolor}{rgb}{.65,.79,.92} % FCUP color

% Configuring the foot line
\setbeamercolor{author in head/foot}{fg=metDarkTeal, bg=barcolor}%
\setbeamercolor{date in head/foot}{fg=barcolor!75!black}%
\setbeamertemplate{footline}
{
  \leavevmode%
  \hbox{%
  \begin{beamercolorbox}[wd=.4\paperwidth,ht=2.25ex,dp=1ex,center]{author in head/foot}%
    \usebeamerfont{author in head/foot}\insertshortauthor
  \end{beamercolorbox}%
  \begin{beamercolorbox}[wd=.5\paperwidth,ht=2.25ex,dp=1ex,center]{title in head/foot}%
    \usebeamerfont{title in head/foot}\insertsection
  \end{beamercolorbox}%
  \begin{beamercolorbox}[wd=.1\paperwidth,ht=2.25ex,dp=1ex,right]{date in head/foot}%
    \insertframenumber{} / \inserttotalframenumber\hspace*{2ex} 
  \end{beamercolorbox}}%
  \vskip0pt%
}
% No configuration symbols
\setbeamertemplate{navigation symbols}{}


%%%%%%% Custom bar above

\setbeamercolor{frametitle}{bg=barcolor,fg=metDarkTeal}
\setbeamercolor{progress bar}{fg=barcolor!75!black}

%%% Custom bar with a LOGO

\makeatletter
\setbeamertemplate{frametitle}{%
  \nointerlineskip%
  \begin{beamercolorbox}[%
      wd=\paperwidth,%
      sep=0pt,%
      leftskip=\metropolis@frametitle@padding,%
      rightskip=\metropolis@frametitle@padding,%
    ]{frametitle}%
  \metropolis@frametitlestrut@start%
  \insertframetitle%
  \nolinebreak%
  \metropolis@frametitlestrut@end%
  \hfill
  \raisebox{-1.3ex}{\includegraphics[height=4ex,keepaspectratio]{images/fcup-bar}}
  \end{beamercolorbox}%
}
\makeatother

%----------------------------------------------------------------------------

\begin{document}

% \setLectureBasic{Cyber-Physical Computation}
\setLecture{2}{Algorithm Correctness}
\frame[plain]{\titlepage}


\section{Motivation}

\byothers{Pedro Ribeiro, slides 1}{1-5}{ribeiro/1_invariants_18092018.pdf}

\section{Specification}

\begin{frame}[fragile]{When is an algorithm correct?}
~\\[-8mm]
\doSimpleExercise{What do these functions do?}
\begin{columns}
\begin{column}{0.48\textwidth}
\begin{lstlisting}[language=C++]
int fa (int x, int y){
  // pre: True
  ...
  // pos: (m == x || m == y) &&
  //      (m >= x && m >= y)
  return m;
}
\end{lstlisting}
\begin{lstlisting}[language=C++]
int fb (int x, int y){
  // pre: x >= 0 && y >= 0
  ...
  // pos: x % r == 0 && y % r == 0
  return r;
}
\end{lstlisting}
\end{column}
\begin{column}{0.48\textwidth}
\begin{lstlisting}[language=C++]
int fc (int x, int y){
  // pre: x > 0 && y > 0
  ...
  // pos: r % x == 0 && r % y == 0
  return r;
}
\end{lstlisting}
\begin{lstlisting}[language=C++]
int fd (int a[], int N){
  // pre: N>0
  ...
  // pos:
  //  (forall_{0<=i<N} x<=a[i]) &&
  //  (exists_{0<=i<N} x==a[i])
  return x;
}
\end{lstlisting}
\end{column}
\end{columns}
\end{frame}

% fa: max
% fb: common divisor
% fc: common multiplier
% fd (ff): max list

% int fd (int a[], int N){
%   // pre: N>0
%   ...
%   // pos: 0 <= p< N &&
%   // forall_{0 <= i< N} a[p] <= a[i]
%   return p;
% }


\begin{frame}{When is an algorithm correct?}
  \doExercise{Formulate pre- and post-conditions:}{
\!\texttt{int prod (int x, int y)} -- product of two integers
\\\texttt{int gcd (int x, int y)} -- greatest common divisor of 2 positive integers
\\\texttt{int sum (int v[], int N)} -- sum of elements in an array
\\\texttt{int maxPOrd (int v[], int N)} -- length of the longest sorted prefix of an array
% its a prefix and there is no next-larger
% forall_{0<i<r} v[i-1]<=v[i])  &&  (r==N || v[r-1]>v[r])
\\\texttt{int isSorted (int v[], int N)} --  tests if an array is sorted (growing)
  }
\end{frame}

\begin{frame}{Hoare triples}
  A triple $\{P\} S \{Q\}$ is a valid Hoare triple when
  \mycbox{if $[P~holds]$ and $[S~is~executed]$ then $[Q~holds]$}

  ~\\[5mm]

  \doExercise{Find initial states that show these are not valid (and fix pre-cond.)}{
  ~\\[-5mm]\begin{enumerate}
    \item \texttt{\{True\} r=x+y; \{r$\geq$x\}} % r>=x[r/x+y] = x+y>=x (y>=0)
    \item \texttt{\{True\} x=x+y; y=x-y; x=x-y; \{x==y\}}
        % going bwd: {x-y == x (x==2y)}; {x==2(x-y) (x==2y)}; {x+y==2y (x==y)}}
    \item \texttt{\{True\} x=x+y; y=x-y; x=x-y; \{x$\neq$y\}}
    \item \texttt{\{True\} if(x>y) r=x-y; else r=y-x; \{r>0\}}
        % opt 1: y-x > 0 && x<=y  ~~~ y>x
        % opt 2: x-y > 0 && x>y   ~~~ x>y
        % hence opt1 || opt2  ~~~ x!=y
    \item \texttt{\{True\} while (x>0) \{y=y+1; x=x-1;\} \{y>x\}}
       % (x>0 && y+x>0 (y>0)) || x<=0 ) -- y=10,x=1
       % x>0 ---> y'=y+x && x'=0
       %     --->  {y+x>0} y=y+x; {y>0} x=0 {y>x}
  \end{enumerate}
  %%%% y>x
  % opt1: x<=0 && y>x
  % opt2: y+1>x-1 && x=1  ~~ y>-1 && x=1
  % opt3: y+2>x-2 && x=2  ~~ y>-2 && x=2
  % opt4: y+3>x-3 && x=3  ~~ y>-3 && x=3 --- y>-x && x>0 ~~~ x>0
  % hence: op1 || x>0  (note: y>0 is sufficient, but not the most generic)
  }
  % fix pre

\end{frame}


\section{Partial correctness}


\begin{frame}{Using rules for Hoare triples}
  $$\begin{array}{c}
    P \Rightarrow Q[x\backslash E]
    \\\hline
    \{P\}~x\texttt{:=}E~\{Q\}
    \end{array}
    \qquad\qquad
  \begin{array}{c}
    P\Rightarrow \alert{I} \qquad \{\alert{I}\land c\} ~S~ \{\alert{I}\} \qquad (\alert{I} \land \lnot c)\Rightarrow Q
    \\\hline
    \{P\} \texttt{ while }c~S~\{Q\}
    \end{array}$$

    \begin{enumerate}
      \item \textbf{Initialisation:} $P \Rightarrow \alert{I}$
        \hfill\textcolor{black!40}{\emph{($P$ is the precondition right before the cycle)}}
        \\Before the cycle the invariant holds.
      \item \textbf{Maintenance:} $\{\alert{I} \land c\}~S~\{ \alert{I} \}$
        \hfill\textcolor{black!40}{\emph{(or $I\land c \Rightarrow I'$, where $I'$ is the invariant after $S$)}}
        \\Assuming the invariant holds before an iteration; it must be valid after it.
      \item \textbf{Usefulness \emph{(Termination)}:} $(\alert{I} \land \lnot c) \Rightarrow Q$
        \hfill\textcolor{black!40}{\emph{(simplify $I\land c$ until obtain $Q$)}}
        \\After the cycle the post-condition holds.
    \end{enumerate}
\end{frame}

\byothers{Pedro Ribeiro, slides 1}{6-10}{ribeiro/1_invariants_18092018.pdf}

\begin{frame}[t]\frametitle{How to show that an invariant is really one?}

  \noindent
  \splitthree{0.30}{0.36}{0.32}{
    \begin{exampleblock}{Initialization}
      The invariant is true prior to the first iteration of the loop
      \\~\\~\\[4mm]
      $P \Rightarrow \alert{I}$
      \\[3mm]\emph{in the slide before:}
      \\$Q \Rightarrow \alert{P}$
    \end{exampleblock}
  }{
    \begin{block}{Maintenance}
      If it is true before an iteration of the loop, it remains true before the next iteration
      \\~\\~\\[4mm]
      $I\land c \Rightarrow \alert{I'}$
      \\[3mm]\emph{in the slide before:}
      \\$P\land B \Rightarrow \alert{P}$ after executing $S$
    \end{block}
  }{
    \begin{alertblock}{Usefulness \emph{(termination)}}
      When the loop terminates, the invariant gives us a useful property that helps show that the algorithm is correct
      \\[4mm]
      $(I \land \lnot c) \Rightarrow \alert{Q}$
      \\[3mm]\emph{in the slide before:}
      \\$P\land \lnot B \Rightarrow \alert{R}$
    \end{alertblock}
  }
\end{frame}


\begin{frame}[fragile]{Exercises}
  
\begin{columns}
\begin{column}{0.48\textwidth}
\begin{lstlisting}[language=C++]
int mult1 (int x, int y){
  // pre: x>=0
  int a, b, r;
  a=x; b=y; r=0;
  while (a!=0){
    r = r+b;
    a = a-1;
  }
  // pos: r == x * y 
  return r;
}
\end{lstlisting}
\end{column}
\begin{column}{0.48\textwidth}
\begin{lstlisting}[language=C++]
int mult2 (int x, int y){
  // pre: x>=0
  int a, b, r;
  a=x; b=y; r=0;
  while (a!=0) {
    if (a%2 == 1) r = r+b;
    a=a/2;
    b=b*2;
  // pos: r == x * y
  return r;
}
\end{lstlisting}
\end{column}
\end{columns}

\only<1>{\doExercise{Check if \emph{Initialization} and \emph{Maintenance} holds for these formulae}{
  ~\\[-5mm]\begin{columns}
    \col[0.2]{$r==a*b$\\$a\geq 0$\\$b\geq 0$}
    \col[0.2]{$r\geq 0$\\$a==x$\\$a\neq x$}
    \col[0.2]{$b== 0$\\$a*b==x*y$\\$a*b+r==x*y$}
  \end{columns}
}}%
\only<2>{\doSimpleExercise{Find loop invariants to prove partial correctness}}
\end{frame}
\exerciseBack

\begin{frame}[fragile]\frametitle{Some intuition -- mult1(4,5)}

\begin{columns}
\begin{column}{0.36\textwidth}
\begin{lstlisting}[language=C++,numbers=left]
int mult1 (int x, int y){
  // pre: x>=0
  int a, b, r;
  a=x; b=y; r=0;
  while (a>0){
    r = r+b;
    a = a-1;
  }
  // pos: r == x * y 
  return r;
}
\end{lstlisting}
\end{column}
\begin{column}{0.25\textwidth}
\footnotesize
\begin{tabular}{c@{~~~~}ccccc}
\toprule
\textbf{line} & \textbf{\texttt{x}} & \textbf{\texttt{y}} &
  \textbf{\texttt{a}} & \textbf{\texttt{b}} & \textbf{\texttt{r}} 
\\\midrule
\textcolor{black!50}{4} & 4 & 5 & 4 & 5 & 0 \\
\midrule
\textcolor{black!50}{6} & 4 & 5 & 4 & 5 & \alert{5} \\
\textcolor{black!50}{7} & 4 & 5 & \alert{3} & 5 & 5 \\
\midrule
\textcolor{black!50}{6} & 4 & 5 & 3 & 5 & \alert{10} \\
\textcolor{black!50}{7} & 4 & 5 & \alert{2} & 5 & 10 \\
\midrule
\textcolor{black!50}{6} & 4 & 5 & 2 & 5 & \alert{15} \\
\textcolor{black!50}{7} & 4 & 5 & \alert{1} & 5 & 15 \\
\midrule
\textcolor{black!50}{6} & 4 & 5 & 1 & 5 & \alert{20} \\
\textcolor{black!50}{7} & 4 & 5 & \alert{0} & 5 & 20 \\
\midrule
\textcolor{black!50}{10} & 4 & 5 & 0 & 5 & 20 \\
\bottomrule
\end{tabular}
\end{column}

\pause
\col[0.32]{
  \begin{itemize}
    \item \texttt{x} and \texttt{y} never change
    \item \texttt{r} grows proportionally as \texttt{a} shrinks
    \item \alert{guess:} $\structure{I} \stackrel{\vartriangle}{~=~} \texttt{a*y} + \texttt{r} = \texttt{x*y}$
    \pause
    \item Need to show:\\[-8mm]
      \begin{align*}
        \texttt{x>=0} &\alert{\Rightarrow} \structure{I'}\\
      %~~~~(where $I' = I[\texttt{x$\backslash$a},\texttt{b$\backslash$y},\texttt{r$\backslash$0}]$)
        \structure{I} \land \texttt{a>0} &\alert{\Rightarrow} \structure{I'}\\
        \structure{I} \land \lnot(\texttt{a>0}) &\alert{\Rightarrow} \texttt{r} = \texttt{x*y}
        ~\\[-12mm]
      \end{align*}
    \pause
    \item (Not all works -- enrich invariant!)
  \end{itemize}
}
\end{columns}
\end{frame}


\begin{frame}[fragile]\frametitle{More exercises}
\begin{columns}
\begin{column}{0.48\textwidth}
% 0,1,2,..,n
% n=0, r=0*...
% n=1, r=1*2/2=1 ()
\begin{lstlisting}[language=C++]
int serie(int n){
  // pre: n>=0
  int r=0, i=1;
  // inv: ??
  while (i!=n+1) {
    r = r+i; i = i+1;
  }
  // pos: r == n * (n+1) / 2;
  return r;
}
\end{lstlisting}
\end{column}
\begin{column}{0.48\textwidth}
\begin{lstlisting}[language=C++]
int mod(int x, int y) {
  // pre: x>=0 && y>0
  int r = x;
  while (y <= r) {
    r = r-y;
  }
  // pos: 0 <= r < y && exists_{q} x == q*y + r
  return r;
}
\end{lstlisting}
\end{column}
\end{columns}

\doSimpleExercise{Find loop invariants}
\end{frame}


\begin{frame}[fragile]\frametitle{Even more exercises (@home)}
\begin{columns}
\begin{column}{0.48\textwidth}
\begin{lstlisting}[language=C++,style=tiny]
int minInd (int v[], int N) {
  // pre: N>0
  int i = 1, r = 0;
  // inv: ???
  while (i<N) {
    if (v[i] < v[r]) r = i;
    i = i+1; }
  // pos: 0 <= r < N && forall_{0 <= k < N} v[r] <= v[k]
return r; }

int minimum (int v[], int N) {
  // pre: N>0
  int i = 1, r = v[0];
  // inv: ???
  while (i!=N) {
     if (v[i] < r) r = v[i];
     i=i+1; }
  // pos: (forall_{0 <= k < N} r <= v[k]) &&
  //      (exists_{0 <= p < N} r == v[p])
  return r;
}

int sum (int v[], int N) {
  // pre: N>0
  int i = 0, r = 0;
  // inv: ???
  while (i!=N) {
    r = r + v[i]; i=i+1;
  }
  // pos: r == sum_{0 <= k < N} v[k]
  return r;
}
\end{lstlisting}
\end{column}
\begin{column}{0.48\textwidth}
\begin{lstlisting}[language=C++,style=tiny]
int sqr1 (int x) {
  // pre: x>=0
  int a = x, b = x, r = 0;
  // inv: ??
  while (a!=0) {
    if (a%2 != 0) r = r + b;
    a=a/2; b=b*2;
  }
  // pos: r == x^2
  return r;
  }

int sqr2 (int x){
  // pre: x>=0
  int r = 0, i = 0, p = 1;
  // inv: ??
  while (i<x) {
    i = i+1; r = r+p; p = p+2;
  }
  // pos: r == x^2
  return r;
}

int ssearch (int x, int a[], int N){
  // pre: N>0 &&
  //      forall_{0 < k < N-1} a[k-1]<=a[k]
  int p = -1, i = 0;
  // inv: ??
  while (p == -1 && i < N && x >= a[i]) {
    if (a[i] == x) p = i;
    i = i+1;
  }
  // pos: (p == -1 && forall_{0 <= k < N} a[k] != x) ||
  //      ( (0 <= p < N ) && x == a[p])
  return p;
}
\end{lstlisting}
\end{column}
\end{columns}
\end{frame}

% inv: (p == -1 && forall_{0 <= k < i} a[k] != x) ||
%      ( (0 <= p < i ) && x == a[p])
% init - ok

% Maint - 
%   IF p == -1 && i < N && x >= a[i] AND a[i] == x, then

%   (p' == -1 && forall_{0 <= k < i'} a[k] != x) ||
%      ( (0 <= p' < i' ) && x == a[p'])
% = (i == -1 && forall_{0 <= k < i+1} a[k] != x) ||
%      ( (0 <= i < i+1 ) && x == a[i]) --> 2nd part true

%   IF p == -1 && i < N && x >= a[i] AND a[i] != x, then

%   (p' == -1 && forall_{0 <= k < i'} a[k] != x) ||
%      ( (0 <= p' < i' ) && x == a[p'])
% = (p == -1 && forall_{0 <= k < i+1} a[k] != x) ||
%      ( (0 <= p < i+1 ) && x == a[p])
% = (p == -1 && forall_{0 <= k < i} a[k] != x && a[i]!=x) -- 1st part true

% Util -
%   IF p != -1 || i >= N || x < a[i]
%   AND (p == -1 && forall_{0 <= k < i} a[k] != x) ||
%      ( (0 <= p < i ) && x == a[p])
%   THEN post?
%     - ( (0 <= p < N ) && x == a[p]) is true when i>=N (&& i<=N inv), p!=1 is false, and x<a[i] is any
%     - (p == -1 && forall_{0 <= k < N} a[k] != x) is true when ... (need pre cond) p!=1 is true and the rest is any... (either i=N and 3rd any, or i<N and a[k]...?)


\section{Complete correctness}

\begin{frame}{Partial/Complete correctness}

~\mycbox[0.6]{Given $\{P\}~S~\{Q\}$}~\\[10mm]

\begin{alertblock}{Partial correctness}
  if $[P~holds]$ and $[S~is~executed]$ then $[Q~holds]$
\end{alertblock}

\begin{alertblock}{Complete correctness}
  if $[P~holds]$ and $[S~is~executed]$ then $[Q~holds]$ \alert{AND $S$ terminates}
\end{alertblock}

\pause
~\mycbox[0.8]{\textbf{Enough to show the existence of a \alert{loop variant}}}

\end{frame}


\begin{frame}[fragile]\frametitle{Loop variant}
  
  Technique that measures the distance between the current state and the final state.

  \begin{block}{A loop variant $\structure{V}$ is an \textbf{integer expression} s.t.}
    \begin{itemize}
      \item is positive in the beginning of each round ($c \land I ~\Rightarrow~ \structure{V}>0$)
      \item decreases in every round ($c \land I ~\Rightarrow~ \structure{V} > \structure{V'}$)
    \end{itemize}
  \end{block}

\begin{columns}
\begin{column}{0.48\textwidth}
\begin{lstlisting}[language=C++]
r=x;
q=0;
while (y <= r) {
  r = r-y;
  q = q+1;
}
\end{lstlisting}
\end{column}
\begin{column}{0.48\textwidth}
  \begin{itemize}
    \item $\structure{V} = r-y$ is not a good variant
    \item \only<1>{...}\only<2->{$\structure{V} =r-y+1$ is a good variant}
      \only<3->{
      \\[3mm] $~~~y\leq r \alert{\Rightarrow} \structure{V} > 0$ at each round
      % V>0 <=> r-y+1 > 0 <=> y<r+1 <=> y<=r (OK)
      \\ $~~~\structure{V} > \structure{V'}$ after each round
    }

  \end{itemize}
\end{column}
\end{columns}

\end{frame}

\begin{frame}[fragile]\frametitle{Exercises}

\begin{lstlisting}[language=C++]
int sum(int v[], int N) {
  int i = 0, r = 0;
  while (i!=N) {
    // variant: ???
    r = r + v[i];
    i=i+1;
  }
  return r;
}
\end{lstlisting}

\doSimpleExercise{Find variant above}
~\\[-10mm]
\doSimpleExercise{Find variants of the loops in previous exercises\\(when searching for invariants)}


\end{frame}

\end{document}
